% \noindent
% \textbf{Setup 2 (Human Evaluation with Real Users)}\quad
Human evaluation is essential for assessing systems designed for collaborative discourse, as it captures the complexities of human interaction, reflects familiar real-world interactions, and provides critical insights into the system’s effectiveness. We compare \system with two baselines: (1) \textbf{RAG Chatbot} as detailed in section \refsec{sec:auto_eval}. (2) Traditional \textbf{Search Engine}.
\subsection{Evaluation Setup}
We conduct an IRB-approved human evaluation to compare \system with RAG Chatbot and Search Engine by recruiting 20 volunteers on the Internet. Participants are randomly split into two groups: one compared \system with Google Search, while the other group compared it with the RAG Chatbot. Participants are asked to seek information on two topics, one on each system, from the same domain and sharing the same goal. Note that we mitigate topic familiarity bias by using two different topics within the same domain.
We have prepared five different domains, with each assigned to 2 users in each group. To counterbalance, one user starts with \system and switches to the baseline for the other topic, and vice versa. 
% \iffalse
% hard to understand:
% We use two different topics to avoid topic familiarity bias.
% We control for confounding variables by creating five pairs of complex information-seeking tasks, each pair consisting of two topics within the same domain and a shared information-seeking goal. Each pair is assigned to two participants per group, with each participant using both systems (one per topic) to avoid topic familiarity bias. We also alternate the order of \system and the baseline to avoid order effects.\fi

%I added the point about topic familiarity. The rest is a repeat. To control for confounding variables, we design five pairs of tasks, where each pair shares a common goal but involves different topics. Each pair is assigned to two participants per group, with each participant interacting with both systems (one per topic) to reduce familiarity bias. Additionally, we alternate the order of using \system and the baseline to minimize order effects.” \yucheng{I add this part back as it describes how we setup human evaluation and how we mitigate confounding variables. I revise the wording to make it easier to understand.}

After seeking information for each topic, participants are instructed to rate their experience based on four grading aspects defined in \refsec{sec:eval_metrics} (\textit{Relevance}, \textit{Breadth}, \textit{Depth}, \textit{Novelty/Serendipity}), using a 5-point Likert scale. After completing both tasks, participants are asked to provide pairwise preferences regarding the \textit{required effort}, \textit{user engagement}, \textit{addressing echo chamber issues}, and \textit{overall experience}. We also collect open-ended feedback and allow participants to optionally leave comments on each discourse turn and the mind map snapshots when interacting with \system. More details on the human evaluation are included in Appendix~\ref{appendix:human_eval}.